% =====================================================================
%  Kollisionskaskaden-HPC -- Projektbericht
%  Paul Kamm, TH Luebeck, HPC via PC2 Paderborn (Noctua 2)
%
%  Bau:   pdflatex report.tex ; pdflatex report.tex   (2x fuer Refs)
%
%  Bilder: PNG only
%    Alte Folien-Plots:   Folien/figures_final/
%    Neue Python-Plots:   results/images/   (je nach Ausfuehrungsordner)
%  Plot-Skripte: scripts/python/ (eigener Ordner)
% =====================================================================
\documentclass[11pt,a4paper]{article}

\usepackage[utf8]{inputenc}
\usepackage[T1]{fontenc}
\usepackage[ngerman]{babel}
\usepackage{lmodern}
\usepackage{microtype}

\usepackage{graphicx}
\usepackage{booktabs}
\usepackage{array}
\usepackage{caption}
\usepackage{subcaption}
\usepackage{amsmath,amssymb}
\usepackage{bm}
\usepackage{siunitx}
\usepackage{geometry}
\usepackage{hyperref}
\usepackage{xcolor}
\usepackage{tcolorbox}
\usepackage{enumitem}
\usepackage{listings}
\usepackage{placeins}   % \FloatBarrier: Grafiken nahe an ihrem Abschnitt halten
\usepackage{float}      % [H]: Grafiken exakt an ihrer Textstelle verankern

\geometry{margin=2.2cm,a4paper}
\raggedbottom   % keine gestreckten Zwischenraeume auf reinen Grafik-Seiten
\hypersetup{colorlinks=true,linkcolor=black,urlcolor=blue!60!black,
            pdftitle={Kollisionskaskaden-HPC Bericht}}
\captionsetup{font=small,labelfont=bf,labelsep=colon}

\graphicspath{%
  {./}%
  {Folien/}%
  {./Folien/}%
  {Folien/figures_final/}%
  {./Folien/figures_final/}%
  {plot_skripte/figures/}%
  {./plot_skripte/figures/}%
  {results/images/}%
  {./results/images/}%
}

% ---- Hilfsmakros -----------------------------------------------------------
\newcommand{\ueberleitung}[1]{\par\smallskip\noindent\textcolor{blue!50!black}{%
  \small $\triangleright$~\emph{#1}}\par\smallskip}

% Infobox für Schlüsselzahlen
\newcommand{\infobox}[2]{%
  \begin{tcolorbox}[colback=orange!5,colframe=orange!60!black,boxrule=0.4pt,
                    arc=2pt,title={\small\textbf{#1}}]
    \small #2
  \end{tcolorbox}%
}

% =====================================================================
\title{Kollisionskaskaden in 2D\\[0.2em]
\large Ein HPC-Projektbericht: GPU-Beschleunigung und Potenzgesetze}
\author{Paul Kamm\\
\small TH L\"ubeck -- Kursprojekt \emph{High Performance Computing}, betreut von Dr.\ Maurice Maurer\\
\small Messungen auf \emph{Noctua 2}, PC2 Paderborn (EPYC 7763, A100-SXM4-40GB)}
\date{\today}

\begin{document}
\maketitle

\begin{abstract}
\noindent
Dieser Bericht untersucht 2D-Kollisionskaskaden aus zwei Perspektiven: als
\textbf{HPC Projekt} und als \textbf{physikalisches Skalierungsgesetz}.
\textbf{Teil~I} (HPC) verfolgt die Frage, wie weit sich eine
Molekulardynamik-Simulation durch Parallelisierung und GPU-Portierung beschleunigen
lässt -- und wo die prinzipielle Grenze liegt. Es zeigt sich, dass der Code
durchgehend \emph{speicherbandbreiten-gebunden} ist: eine A100-GPU liefert
\textbf{731-fache} Beschleunigung gegenüber einem EPYC-Kern; die MPI-Skalierung
ist U-förmig mit superlinearer Effizienz (133\,\% bei 128 Kernen) durch
Cache-Effekte; Kommunikationsoptimierungen bringen nichts.
\textbf{Teil~II} (Power Law) zeigt, dass die Clustergrößenverteilung
aus 1000 simulierten Kaskaden einem Potenzgesetz mit $S = 1{,}38$--$1{,}41$
folgt -- robust gegen Healing. Behandelt werden die Fit-Methodik (Maximum
Likelihood nach Clauset et al., $x_\text{min}$-Falle) sowie die Einordnung des Ergebnisses
in den Kontext der 3D-MD-Befunde von Sand et al.\,(2013) und der
Subkaskaden-Theorie von De Backer et al.\,(2016).
\end{abstract}

\newpage
{\footnotesize
\makeatletter
\renewcommand*\l@part[2]{%
  \ifnum \c@tocdepth >-2\relax
    \addpenalty\@secpenalty
    \addvspace{1.3em \@plus\p@}%
    \setlength\@tempdima{3em}%
    \begingroup
      \parindent \z@ \rightskip \@pnumwidth
      \parfillskip -\@pnumwidth
      {\leavevmode
       \large \bfseries #1\hfil
       \hb@xt@\@pnumwidth{\hss #2%
                          \kern-\p@\kern\p@}}\par
       \nobreak
    \endgroup
  \fi}
\renewcommand*\l@section[2]{%
  \ifnum \c@tocdepth >\z@
    \addpenalty\@secpenalty
    \addvspace{0.4em \@plus\p@}%
    \setlength\@tempdima{1.5em}%
    \begingroup
      \parindent \z@ \rightskip \@pnumwidth
      \parfillskip -\@pnumwidth
      \leavevmode \bfseries
      \advance\leftskip\@tempdima
      \hskip -\leftskip
      #1\nobreak\hfil
      \nobreak\hb@xt@\@pnumwidth{\hss #2%
                                 \kern-\p@\kern\p@}\par
    \endgroup
  \fi}
\makeatother
\tableofcontents}
\vspace{1em}

% =====================================================================
% TEIL I: HPC
% =====================================================================
\newpage
\FloatBarrier
\part{HPC -- Beschleunigung der Simulation}
\label{part:hpc}

% =====================================================================
\FloatBarrier
\section{Das Simulations-Modell}
\label{sec:modell}

\subsection{Physikalisches Modell: 2D-Dreiecks\-gitter}

Simuliert wird kein echtes Material, sondern der \emph{Mechanismus} einer
Kollisionskaskade in einem minimalen 2D-Modell. Die Atome sitzen auf einem
\textbf{Dreiecksgitter} (6er-Koordination, schubstabil, elastisch isotrop;
Abb.~\ref{fig:modell-schema}).

\paragraph{Zwei Wechselwirkungen (Abb.~\ref{fig:kraftgesetz}).}
\begin{itemize}[itemsep=2pt]
\item \textbf{Brechbare harmonische Feder:} $F = K(r - L_0)$ mit Bruch bei
  $r > 1{,}15\,L_0$ ($K = 100$). Modelliert die kovalente/metallische Bindung.
\item \textbf{Steile Kurzreichweitswechselwirkung:} $F = K_\text{rep}
  \bigl((R_\text{cut}/r)^{12} - 1\bigr)$ mit $K_\text{rep} = 400$,
  $R_\text{cut} = 0{,}9\,L_0$. Gibt den Atomen einen effektiven
  Wirkungsquerschnitt -- ohne sie würde das PKA alles durchdringen.
\end{itemize}

Das PKA (Primary Knock-on Atom) ist schlicht ein Atom mit einer hohen
Anfangsgeschwindigkeit $v_\text{PKA} = \sqrt{2E_\text{PKA}/m}$, die in eine
Richtung zeigt. Im Laufe der Simulation \"ubertr\"agt es seine Energie
durch Absto\ss ungen auf Nachbaratome, die ihrerseits weitere ansto\ss en --
die verzweigende Kaskade (Abb.~\ref{fig:pka}).

\begin{figure}[htbp]
  \centering
  \includegraphics[width=0.9\linewidth]{M1_model_schema.png}
  \caption{Modell-Schema. Links: intaktes Dreiecksgitter mit rot markiertem PKA.
    Rechts: Zustand nach der Kaskade -- gerissene Bindungen (rot) und ein
    Feld aus Defektclustern. Alle Größen sind dimensionslos.}
  \label{fig:modell-schema}
\end{figure}

\begin{figure}[htbp]
  \centering
  \includegraphics[width=0.9\linewidth]{M2_force_law.png}
  \caption{Die beiden Wechselwirkungen. Links: Kraftverläufe -- Feder (blau,
    Bruch bei 1{,}15\,$L_0$) und Abstoßung (rot, nur für $r < 0{,}9\,L_0$).
    Rechts: zugehörige Potenziale mit stetigem Nullpunkt an $R_\text{cut}$.
    Die Energiestetigkeit ist Grundlage des Korrektheitschecks.}
  \label{fig:kraftgesetz}
\end{figure}

\begin{figure}[htbp]
  \centering
  \includegraphics[width=0.85\linewidth]{M4_pka_velocity.png}
  \caption{PKA-Start (Korrektheitsprüfung). Links: kinetische Gesamtenergie
    über die Zeit -- bei $t=0$ steckt sie vollständig im PKA-Atom und wird
    im Zuge der Kaskade auf die Nachbaratome verteilt; rechts:
    Geschwindigkeitshistogramm bei $t=0$, ein
    einzelner Peak bei $v_\text{PKA} = \sqrt{2E/m} \approx 100$, alle anderen
    Atome stehen still. Genau dieser Impuls startet die Kaskade.}
  \label{fig:pka}
\end{figure}

\newpage

\subsection{Datenstruktur: Structure of Arrays (SoA)}
\label{sec:soa}

Die Atome werden als \textbf{Structure of Arrays} gespeichert: statt
\texttt{atom[i].px} gibt es separate Arrays \texttt{px[i]}, \texttt{py[i]},
\texttt{vx[i]}, \texttt{vy[i]}, \texttt{fx[i]}, \texttt{fy[i]},
\texttt{nbr[i][6]}, \texttt{cell[i]}.

\infobox{Speicherlayout (76\,B/Atom)}{%
  8 Felder: \texttt{px,py,vx,vy,fx,fy} (je 8\,B $\times$ 6) + \texttt{nbr}
  (4\,B $\times$ 6, sechs Nachbarn im Dreiecksgitter) + \texttt{cell} (4\,B)
  = 76\,B/Atom.
  Bei $N = 2\,\text{Mio}$: Gesamtdatensatz $\approx 152\,$MB.
  Dieses Maß wird in Abschnitt~\ref{sec:cache} zur Cache-Analyse verwendet.
}

SoA hat drei Vorteile für diese Simulation:
\begin{enumerate}[itemsep=2pt]
\item \textbf{Cache-Effizienz:} Der Kraftkern iteriert über $N$ Atome und
  braucht dabei $px$, $py$, $nbr$ sequenziell -- bei SoA ist das ein
  einziger linearer Speicherzugriff.
\item \textbf{Vektorisierbarkeit:} SIMD-Einheiten (AVX-512 auf EPYC)
  können mehrere Atome gleichzeitig laden und rechnen.
\item \textbf{GPU-Kompatibilität:} CUDA-Kernel können mit einem Thread pro
  Atom arbeiten und koaleszente Speicherzugriffe machen (Abschnitt~\ref{sec:cuda}).
\end{enumerate}

\subsection{Ortsraum-Zerlegung: Cell List}
\label{sec:celllist}

Ohne weitere Maßnahme wäre die Komplexität des Kraftkerns $O(N^2)$ -- für
2\,Mio Atome nicht praktikabel. Die \textbf{Cell List} löst das
(Abb.~\ref{fig:celllist}):

\begin{figure}[htbp]
  \centering
  \includegraphics[width=0.45\linewidth]{M3_cell_list.png}
  \caption{Cell-List-Schema. Der Raum wird in Zellen der Größe $R_\text{cut}$
    unterteilt. Für ein Atom (rot) genügt es, die 9 Nachbarzellen (blaues Band)
    zu durchsuchen. Verkettete Listen über \texttt{cell\_head}/\texttt{cell\_next}
    machen die Nachbarsuche $O(N)$ statt $O(N^2)$.}
  \label{fig:celllist}
\end{figure}

Jede Zelle hat Seitenlänge $\geq R_\text{cut}$; potenzielle Bindungspartner
liegen garantiert in den 9 umgebenden Zellen. Für $N$ Atome mit konstanter
Dichte gilt: mittlere Atome pro Zelle $= O(1)$, also Gesamtkomplexität $O(N)$.
Die Cell List muss nach jeder Migration neu aufgebaut werden, was im
MPI-Zeitschritt als eigenständige Phase (\texttt{cells}) gemessen wird.

\subsection{Integrator: Velocity-Verlet}
\label{sec:integrator}

Die Bewegungsgleichungen werden mit \textbf{Velocity-Verlet} in
Kick-Drift-Kick-Form integriert:
\begin{align}
\bm{v}_{i}^{n+1/2} &= \bm{v}_{i}^{n} + \tfrac{\Delta t}{2}\,\bm{F}_{i}^{n}/m
  \quad \text{(Kick)} \label{eq:kick1}\\
\bm{r}_{i}^{n+1}   &= \bm{r}_{i}^{n} + \Delta t\,\bm{v}_{i}^{n+1/2}
  \quad \text{(Drift)} \label{eq:drift}\\
\bm{v}_{i}^{n+1}   &= \bm{v}_{i}^{n+1/2} + \tfrac{\Delta t}{2}\,\bm{F}_{i}^{n+1}/m
  \quad \text{(Kick)} \label{eq:kick2}
\end{align}
Als symplektischer Integrator erhält Velocity-Verlet die Phasenraum-Struktur
und zeigt keinen systematischen Energiedrift -- im Gegensatz zum
expliziten Euler-Verfahren. Abbildung~\ref{fig:energiecheck} belegt, dass die
Gesamtenergie über 3000 Schritte unter $0{,}6\,\%$ (seriell) bzw.\
$0{,}22\,\%$ (A100, fp64) driftet. Abbildung~\ref{fig:timeseries} zeigt die
zugehörige Zeitentwicklung der Kaskade.

\begin{figure}[H]
  \centering
  \includegraphics[width=0.8\linewidth]{vid_S1_energy.png}
  \caption{Energiebilanz (Korrektheitsprüfung). Die Gesamtenergie $E_\text{total}$
    (dicke schwarze Linie) bleibt praktisch konstant. Der prozentuale Drift
    liegt unter $0{,}6\,\%$ für die serielle CPU-Version und unter $0{,}22\,\%$
    für den A100 (fp64). Ohne diesen Nachweis wäre jede Performance-Messung
    wertlos.}
  \label{fig:energiecheck}
\end{figure}

\begin{figure}[H]
  \centering
  \includegraphics[width=0.8\linewidth]{vid_S6_timeseries.png}
  \caption{Kaskaden-Zeitentwicklung ($t = 0 \dots T$). Farbe = Geschwindigkeitsbetrag.
    Ballistische Phase (Ausbreitung der Stoßfront) klingt ab $t \sim 1$ ab;
    danach steht die Defektstruktur. Die Zeitskala $T = 3000\,\Delta t$ ist
    gro\ss\  genug, damit die Kaskade vollständig zur Ruhe kommt.}
  \label{fig:timeseries}
\end{figure}

% =====================================================================
\FloatBarrier
\section{MPI-Parallelisierung}
\label{sec:mpi}

\subsection{Zerlegungs-Strategie}

\paragraph{1D-Domänenzerlegung in $y$-Richtung.}
Der 2D-Raum wird in horizontale Streifen zerlegt, jeweils ein Streifen pro
MPI-Rank. Die 1D-Zerlegung hat einen entscheidenden Vorteil: von den $P-1$
Rank-Grenzen kreuzt beim Netzwerkwechsel (z.B. von Knoten 1 zu Knoten 2)
genau eine einzige -- die Inter-Node-Kommunikation bleibt minimal, wie
Abschnitt~\ref{sec:nodecomp} belegt.

\paragraph{Gather-Formulierung.}
Jedes Atom summiert \emph{nur seine eigene Kraft}, Newton 3 wird nicht
genutzt. Das bedeutet: jedes Atompaar $(i,j)$ wird zweimal ausgewertet.
Diese scheinbar ineffiziente Formulierung hat drei Gründe:
\begin{enumerate}[itemsep=2pt]
\item \textbf{Korrektheit über Rank-Grenzen}: Ein Geisteratom (Kopie aus dem
  Nachbar-Rank) hat keinen ``Besitzer'', der die Hälfte der Kraft
  einsammeln würde.
\item \textbf{Atomicity-Freiheit auf der GPU}: Kein \texttt{atomicAdd} nötig
  (Abschnitt~\ref{sec:cuda}).
\item \textbf{Cache-Effizienz}: Jeder Thread schreibt nur an \emph{einen}
  Speicherplatz.
\end{enumerate}

\subsection{Phasen im MPI-Zeitschritt}

Jeder MPI-Zeitschritt besteht aus sechs messbaren Phasen:

\begin{center}\small
\begin{tabular}{@{}lllr@{}}
\toprule
Phase & Aufgabe & Typ & Anteil @\,1 Kern \\
\midrule
\texttt{force}   & Federn + Abstoßung + Bindungsbruch (Gather)   & Rechnung   & 89\,\% \\
\texttt{hash}    & gid $\to$ lokaler Index (für Bindungsbruch)    & Rechnung   &  6\,\% \\
\texttt{cells}   & Cell List neu aufbauen (lokal + Halo)          & Rechnung   &  2\,\% \\
\texttt{halo}    & Geisteratome von Nachbar-Ranks kopieren        & Komm.      & $<$1\,\% \\
\texttt{migrate} & Besitzwechsel beim Grenzübertritt              & Komm.      & $<$1\,\% \\
\texttt{energy}  & Globale Energiediagnose (\texttt{MPI\_Reduce}) & Diagnose   & $<$1\,\% \\
\bottomrule
\end{tabular}
\end{center}

\textbf{Halo-Austausch.}
Jeder Rank sendet eine Kopie seiner Randatome (``Geisteratome'') an die
beiden Nachbar-Ranks. Das passiert jeden Zeitschritt, weil die Federn
über Ranggrenzen hinweg wirken können. Das Halo-Volumen ist proportional
zur Gitterbreite $N_x$ (konstant bei 1D-Zerlegung), was den Kommunikationsanteil
bei kleinen Rank-Zahlen harmlos hält, aber bei vielen Ranks relativ wächst.

\textbf{Atom-Migration.}
Überquert ein Atom seine Rank-Grenze, wechselt es den ``Besitzer''. Das
Atom wird mit allen Daten (Position, Geschwindigkeit, Bindungstabelle)
an den neuen Rank gesendet und beim alten gelöscht. In einer Kaskade
betrifft das typischerweise nur wenige Atome pro Schritt.

\subsection{Speicherkorrektheit: Valgrind-Analyse}
\label{sec:valgrind}

Neben der physikalischen Korrektheit (Energieerhaltung, Abb.~\ref{fig:energiecheck})
prüft \emph{Valgrind Memcheck} zusätzlich die Speicherkorrektheit beider
CPU-Implementierungen: ungültige Zugriffe, Verwendung uninitialisierten
Speichers und Speicherlecks. Da Memcheck jeden Speicherzugriff
instrumentiert und dadurch $20$--$50\times$ langsamer läuft, wurde nicht das
$2$-Mio.-Atome-Benchmarkgitter verwendet, sondern ein kleines Gitter
($80\times80 = 6400$ Atome, $300$ Schritte, \texttt{params\_valgrind.ini}).
Dieses durchläuft dieselben Code-Pfade (Allokation, Kraftberechnung, Cell-List,
Halo-Austausch, Migration), ist aber in Sekunden statt Stunden analysiert.

\begin{lstlisting}[basicstyle=\ttfamily\footnotesize, breaklines=true]
# Serielle Version (mit gcc gebaut, damit keine MPI-Bibliothek stoert):
gcc -O1 -g -lm -march=native -o cascade_serial cascade_serial.c
valgrind --leak-check=full --show-leak-kinds=all --track-origins=yes \
         ./cascade_serial params_valgrind.ini

# MPI-Version (2 Raenge, ein Logfile pro Rang):
mpicc -O1 -g -lm -march=native -o mpi_cascade mpi_cascade.c
mpirun --oversubscribe -np 2 valgrind --leak-check=full \
       --show-leak-kinds=definite,indirect --track-origins=yes \
       --log-file=vg_rank_%q{OMPI_COMM_WORLD_RANK}.txt \
       ./mpi_cascade params_valgrind.ini
\end{lstlisting}

\begin{center}\small
\begin{tabular}{@{}lrr@{}}
\toprule
Kategorie & \texttt{cascade\_serial} & \texttt{mpi\_cascade} (je Rang) \\
\midrule
Ungültige Zugriffe (App) & 0 & 0 \\
Uninit.\ Werte (App)     & 0 & 0 \\
\emph{definitely lost} (App) & 0 & $\approx 1{,}22$\,MB \\
\emph{still reachable} (App) & $\approx 0{,}96$\,MB & --- \\
Fehler in MPI-Bibliothek & --- & $\sim 10^7$ (Fehlalarme) \\
\bottomrule
\end{tabular}
\end{center}

\paragraph{Serielle Version: sauber.} Memcheck meldet
\emph{0 errors from 0 contexts}. Die $\approx 0{,}96$\,MB \emph{still reachable}
sind die einmalig allokierten SoA-Arrays, die bewusst bis Programmende gehalten
werden; alle Allokationen pro Zeitschritt werden korrekt freigegeben.

\paragraph{MPI-Version: keine echten Fehler im Anwendungscode.} Die $\sim 10^7$
gemeldeten Fehler stammen \emph{sämtlich} aus dem Inneren von OpenMPI
(\texttt{MPI\_Init}) -- ein bekannter, harmloser Valgrind-Fehlalarm; kein
einziger Fehler-Kontext zeigt auf \texttt{mpi\_cascade.c}. Die
$\approx 1{,}22$\,MB \emph{definitely lost} sind die einmalig zum Start
allokierten Strip-SoA-Felder, deren Zeiger beim Leak-Check (nach Rückkehr aus
\texttt{main}) nicht mehr erreichbar sind. Es ist \emph{kein} mit der Laufzeit
wachsendes Leck: den $44\,386$ Allokationen stehen $43\,730$ Freigaben
gegenüber -- die Differenz sind genau die Startup-Felder.

\infobox{Schlüsselergebnis Valgrind}{%
Beide CPU-Versionen sind frei von ungültigen Speicherzugriffen und
Verwendung uninitialisierter Werte im Anwendungscode. Es gibt \emph{kein}
mit der Laufzeit wachsendes Speicherleck; sämtliche nicht freigegebene
Speicher sind einmalige Startup-Allokationen, die das Betriebssystem beim
Programmende zurückgibt. Die Fehlermeldungen der MPI-Version sind
Bibliotheks-interne Valgrind-Fehlalarme aus \texttt{MPI\_Init}.}

% =====================================================================
\FloatBarrier
\section{Compiler-Optimierungen (Einzelkern)}
\label{sec:compiler}

Vor der Parallelisierung stellt sich die Frage: Wie viel bringt der Compiler?
Gemessen mit verschiedenen Flag-Kombinationen auf einem EPYC-7763-Kern.

\infobox{Schlüsselergebnis Compiler}{%
  Baseline \texttt{-O0}: 113{,}7\,s. Schnellstes \emph{korrektes} Setup:
  \texttt{-O3} mit 57{,}0\,s ($\approx 2{,}0\times$). \texttt{-Ofast/-ffast-math}
  ist \emph{langsamer} und verletzt die Energieerhaltung.
}

\begin{figure}[htbp]
  \centering
  \includegraphics[width=0.85\linewidth]{H9_compiler_opt.png}
  \caption{Compiler-Optimierungen auf einem EPYC-7763-Kern (Noctua 2), sortiert
    nach Wandzeit. Grün markiert ist \texttt{-O3} als empfohlenes Setup.
    \texttt{-O3} liefert $\approx 2\times$ Speedup gegenüber \texttt{-O0}.
    \texttt{-Ofast} und \texttt{-ffast-math} umgehen IEEE-754-Regeln und
    erzeugen nicht-assoziative Floating-Point-Operationen, was in einem
    symplektischen Integrator grundsätzlich riskant für die Energieerhaltung ist.}
  \label{fig:compiler}
\end{figure}

Warum bringt \texttt{-O3} nur den Faktor 2? Die \texttt{force}-Phase ist
\emph{speicherbandbreiten-gebunden}: sie streamt $\sim 152$\,MB Atom-Daten
pro Schritt und rechnet dabei pro Atom nur wenige Floating-Point-Operationen.
Bessere Instruktionen helfen nicht, wenn die Daten nicht schnell genug ankommen.
Dieses Muster zieht sich wie ein roter Faden durch alle folgenden Messungen.

% =====================================================================
\FloatBarrier
\section{Strong Scaling -- fester Problemumfang}
\label{sec:strong}

\textbf{Setup:} $N = 1\,999\,396$ Atome ($1414 \times 1414$), 1--256 Kerne,
Phasentimer aktiv, ein Lauf je Rankzahl (Wiederholungsmessungen mit
identischem Setup liefert der Optimierungsvergleich in Abschnitt~\ref{sec:opt}).

\begin{figure}[htbp]
  \centering
  \includegraphics[width=\linewidth]{H1_strong_scaling.png}
  \caption{Strong Scaling. Oben: Speedup $S = T(1)/T(P)$ vs.\ $P$, mit
    idealer Gerade (grau). Unten: Effizienz $S/P$ in \%. Die U-Form mit
    Minimum bei 16 Kernen und superlinearem Maximum ($133\,\%$) bei 128 Kernen
    ist kein Messfehler -- Abschnitt~\ref{sec:cache} erklärt sie quantitativ
    durch Cache-Effekte.}
  \label{fig:strong}
\end{figure}

\paragraph{Schlüsselzahlen.}
\begin{center}\small
\begin{tabular}{@{}rrrr@{}}
\toprule
Kerne & $T(P)$ [s] & Speedup & Effizienz \\
\midrule
1   & 1152{,}9 & 1{,}0   & 100\,\% \\
4   &  360{,}1 & 3{,}2   &  80\,\% \\
8   &  217{,}2 & 5{,}3   &  66\,\% \\
16  &  152{,}6 & 7{,}6   &  47\,\% (Minimum) \\
32  &   59{,}5 & 19{,}4  &  61\,\% \\
64  &   22{,}5 & 51{,}2  &  80\,\% \\
128 &    6{,}7 & 170{,}8 & 133\,\% (superlinear!) \\
\bottomrule
\end{tabular}
\end{center}

\paragraph{Drei Regionen der Effizienzkurve.}
\begin{itemize}[itemsep=2pt]
\item \textbf{1$\to$16 Kerne (Abstieg):}
  Die 1D-Zerlegung fügt bei jedem neuen Rank eine Halo-Grenze hinzu.
  Das Halo-Volumen pro Rank bleibt konstant (volle Breite $N_x$), während
  die Rechenarbeit pro Rank schrumpft. Bei 16 Kernen liegt der
  Kommunikationsanteil (Maximum über die Ranks) bei $\approx 35\,\%$ der
  Wandzeit; er wächst bis 32 Kerne weiter auf $\approx 60\,\%$
  (Abb.~\ref{fig:phases}), wird dort aber bereits von den einsetzenden
  Cache-Effekten überkompensiert.
\item \textbf{16$\to$64 Kerne (Erholung):}
  Cache-Effekte setzen ein: die Datenmenge pro Rank sinkt unter
  $\approx 9$\,MB und nähert sich dem L3-Cache (4\,MB/Rank).
\item \textbf{64$\to$128 Kerne (superlinear):}
  Datensatz/Rank liegt weit unterhalb des L3-Limits -- der
  \texttt{force}-Kernel wird cache-resident. Der Speedup übersteigt das
  Ideal. Das ist quantitativ vorhergesagt durch das Cache-Argument
  (Abschnitt~\ref{sec:cache}).
\end{itemize}

% =====================================================================
\FloatBarrier
\section{Cache-Argument: Warum die Effizienz superlinear wird}
\label{sec:cache}

\infobox{AMD EPYC 7763 L3-Cache}{%
  AMD EPYC 7763: 8 Cache Dies (CCDs), je 32\,MB L3, je 8 Kerne.
  $\Rightarrow$ L3-Anteil pro Rank: $32\,\text{MB} / 8 = 4\,\text{MB}$.
  SoA-Datensatz: $76\,\text{B/Atom} \times N / P$.
  Bei $P = 64$ Ranks: $76 \times 1{,}999 \cdot 10^6 / 64 \approx 2{,}4$\,MB
  $< 4$\,MB $\Rightarrow$ cache-resident.
}

\begin{figure}[htbp]
  \centering
  \includegraphics[width=0.9\linewidth]{cache_argument.png}
  \caption{Cache-Argument (Dual-Achsen-Plot). Blaue Kurve (linke Achse, log):
    Datensatz pro Rank in MB, sinkt mit $1/P$. Gestrichelte Linie: L3-Schwelle
    bei 4\,MB/Rank. Rote Kurve (rechte Achse): parallele Effizienz. Der
    Datensatz unterschreitet die L3-Schwelle zwischen 32 und 64 Ranks;
    die Effizienz erholt sich dort und übersteigt bei 128 Ranks 100\,\%.
    Skript: \texttt{scripts/python/plot\_cache\_argument.py}.}
  \label{fig:cache}
\end{figure}

Der Mechanismus ist direkt: Die \texttt{force}-Phase streamt die Arrays
\texttt{px, py, nbr} für jedes Atom. Solange der Datensatz nicht in den
L3-Cache passt, muss er bei jedem Schritt aus dem DRAM geladen werden
(Bandbreite: ca.\ 400\,GB/s pro EPYC-Knoten, geteilt unter allen Kernen).
Sobald er in den L3 passt, ist er cache-resident: der Kernel sieht
$\sim 2$\,TB/s effektive Bandbreite -- also $5\times$ so viel. Genau das
erklärt den superlinearen Sprung.

\begin{figure}[htbp]
  \centering
  \includegraphics[width=\linewidth]{H2_phase_breakdown.png}
  \caption{Phasenbreakdown über alle Kerne. Links: absolute Zeiten (log-$y$)
    -- \texttt{force} (blau) fällt überproportional. Rechts: normierter
    Stack -- der Kommunikationsanteil (\texttt{halo}+\texttt{migrate},
    Rottöne) wächst mit der Rankzahl deutlich an, während \texttt{force}
    dominant bleibt.}
  \label{fig:phases}
\end{figure}

\begin{figure}[htbp]
  \centering
  \includegraphics[width=\linewidth]{H3_compute_vs_comm.png}
  \caption{Rechnen vs.\ Kommunikation. Rechenzeit (force+cells+hash, blau)
    fällt schneller als $1/P$ -- superlinear wegen Cache.
    Kommunikationszeit (halo+migrate, rot) wächst zunächst und sinkt erst
    bei voller Knotenbelegung. Der Flaschenhals wandert mit der Rankzahl.}
  \label{fig:compute-comm}
\end{figure}

% =====================================================================
\FloatBarrier
\section{1-Knoten vs.\ 2-Knoten-Vergleich}
\label{sec:nodecomp}

\begin{figure}[htbp]
  \centering
  \includegraphics[width=0.42\linewidth]{H8_node_comparison.png}
  \caption{1-Knoten vs.\ 2-Knoten. Bei überlappenden Rank-Zahlen ($P = 64, 128$)
    sind die Wandzeiten praktisch identisch. In der 1D-Zerlegung
    kreuzt genau eine von $P-1$ Rank-Grenzen den Inter-Node-Link --
    die Netzwerkkosten sind vernachlässigbar.}
  \label{fig:nodes}
\end{figure}

Der 256-Kerne-Punkt (nur 2-Knoten erreichbar) zeigt, dass die Superlinearität
auch über Knoten-Grenzen hinweg anhält. Die 1D-Zerlegung ist robust gegenüber
der Netzwerktopologie: von den $P-1$ Rank-Grenzen sind $(P-2)$ intra-node und
genau eine inter-node.

% =====================================================================
\FloatBarrier
\section{Weak Scaling -- variable Problemgröße}
\label{sec:weak}

\textbf{Setup:} Problemgröße wächst proportional zur Rankzahl (konstante
Atomzahl pro Rank). Im Idealfall konstante Wandzeit.

\begin{figure}[htbp]
  \centering
  \includegraphics[width=\linewidth]{H5_weak_scaling.png}
  \caption{Weak Scaling. Links: Wandzeit -- bis 16 Kerne nahezu ideal,
    dann monoton steigend auf $\approx 162$\,s bei 256 Kernen.
    Rechts: Weak-Effizienz $T_0/T_P$ -- sinkt auf 23\,\%.}
  \label{fig:weak}
\end{figure}

\newpage

\paragraph{Mechanismus: gesättigte Speicherbandbreite.}
Im Weak Scaling bleibt die Datenmenge \emph{pro Rank} konstant -- aber die
\emph{gemeinsame} Speicherbandbreite eines Knotens ist physikalisch fest und
wird von allen Kernen geteilt. Bis 16 Kerne ist sie nicht ausgelastet
(flache Kurve). Ab $\approx 32$ Kernen ist sie gesättigt -- mehr Kerne
liefern mehr Rechenwerke, aber nicht mehr Speicherdurchsatz. Jeder
weitere Kern wartet auf Daten.

\infobox{Einheitliche Diagnose}{%
  Strong Scaling: superlinear wegen Cache. Weak Scaling: bricht ein wegen
  Speicherbandbreite. \textbf{Beide Befunde zeigen denselben Engpass aus
  verschiedenen Blickwinkeln}: der Code ist memory-bound, nicht
  compute-bound und nicht kommunikationsgebunden.
}

% =====================================================================
\FloatBarrier
\section{Optimierungsvergleich: Non-Blocking Halo und Re-Sorting}
\label{sec:opt}

Zwei klassische HPC-Optimierungen wurden als Schalter implementiert und
gemessen:

\begin{itemize}[itemsep=2pt]
\item \textbf{Non-blocking Halo:} \texttt{MPI\_Isend/Irecv} mit
  Richtungs-Overlap -- beide Austauschrichtungen laufen gleichzeitig
  statt nacheinander (kein Überlappen mit Rechnung).
\item \textbf{Periodisches Re-Sorting:} Umsortieren der SoA-Arrays
  nach räumlicher Lokalität -- Cache-Hygiene nach Atom-Migration.
\end{itemize}

\begin{figure}[htbp]
  \centering
  \includegraphics[width=\linewidth]{H6_opt_compare.png}
  \caption{Optimierungsvergleich über 1--128 Kerne. Links: Speedup-Kurven ohne
    (blau) und mit Optimierung (rot) liegen praktisch deckungsgleich. Rechts:
    relative Zeitersparnis liegt im $\pm 2\,\%$-Rauschband, im Mittel
    $+0{,}8\,\%$.}
  \label{fig:opt}
\end{figure}

\paragraph{Warum beide nichts bringen.}
\begin{itemize}[itemsep=2pt]
\item \textbf{Non-blocking Halo halbiert bestenfalls die Zeit der beiden
  Austauschrichtungen.} Das ausgetauschte Datenvolumen ist klein; die
  gemessene \texttt{halo}-Zeit ist weit größer, als das reine Volumen
  erklären würde -- sie ist überwiegend Warte-/Synchronisationszeit auf
  den Nachbar-Rank, und dieses Warten verschwindet durch
  \texttt{Isend/Irecv} nicht.
\item \textbf{Re-Sorting greift nur bei starker Cache-Degradation durch Migration.}
  In einer Kaskade ist die Bewegung räumlich lokalisiert -- kaum ein Atom
  überquert eine Rank-Grenze. Die Speicherordnung degradiert kaum; Re-Sorting
  kostet fast genauso viel wie es spart.
\end{itemize}

\emph{Diese Null-Ergebnisse sind das eigentliche Ergebnis, kein Versagen.}
Sie beweisen unabhängig von den Scaling-Kurven: der Flaschenhals ist die
Speicherbandbreite.

% =====================================================================
\FloatBarrier
\section{CUDA-Portierung}
\label{sec:cuda}

\subsection{Architektur-Entscheidungen}

Die GPU-Version portiert dieselbe Gather-Formulierung auf CUDA:
\textbf{ein Thread pro Atom}, $i = \texttt{blockIdx.x} \cdot \texttt{blockDim.x}
+ \texttt{threadIdx.x}$.

\paragraph{Keine Atomics im Kraftkern.}
Da jeder Thread nur die Kraft auf \emph{sein} Atom berechnet, gibt es keinen
konkurrierenden Schreibzugriff. Jedes Paar $(i,j)$ wird zweimal ausgewertet
(wie in der CPU-Version), aber dafür keine \texttt{atomicAdd}-Operationen.
Atomics sind auf GPU sehr teuer; die doppelte Rechenarbeit ist billiger als
eine Synchronisation auf dem gemeinsamen Datenpfad.

\paragraph{Block-Größe 256.}
Aus dem Sweep über Block-Größen (128, 256, 512, 1024) ergibt sich 512 als
leicht bestes und 256 als praktisch gleichwertig. Verwendet wird $T = 256$
(gut teilbar, gute Auslastung bei $N = 2$\,Mio Atomen: $7811$ Blöcke).

\subsection{Benchmark-Ergebnis: A100-SXM4-40GB}
\label{sec:a100bench}

\infobox{A100 (fp64) Headline-Ergebnis}{%
  $N = 1\,999\,396$ Atome, 3000 Schritte, 5 Wiederholungen.
  Median: \textbf{1{,}5765\,s} = \textbf{525\,µs/Schritt} =
  \textbf{3805\,Matom/s}.\\
  Force-Kernel: 259\,µs/Aufruf (48\,\% der Zeitschleife).\\
  Energie-Drift: 0{,}22\,\% (fp64, symplektisch).\\
  H2D-Transfer (184\,MB): 14\,ms (14\,GB/s).
}

\begin{figure}[htbp]
  \centering
  \includegraphics[width=0.9\linewidth]{tts.png}
  \caption{Time-to-Solution-Vergleich ($N \approx 2$\,Mio Atome, 3000 Schritte).
    Links: Wandzeit auf Log-Skala; Rechts: Speedup relativ zu 1 CPU-Kern.
    Der A100 (fp64) ist \textbf{731-fach} schneller als 1 EPYC-Kern und
    \textbf{14-fach} schneller als 64 EPYC-Kerne.
    Skript: \texttt{scripts/python/plot\_tts.py}.}
  \label{fig:tts}
\end{figure}

\paragraph{fp32 vs.\ fp64.}
Bei fp32 halbiert sich die zu bewegende Datenmenge -- ein
bandbreitengebundener Kernel kann also theoretisch $2\times$ schneller laufen.
Resultat: fp32 läuft in 0{,}999\,s ($\approx 1{,}6\times$ schneller als fp64),
aber der Energie-Drift steigt auf $0{,}94\,\%$ -- noch tolerierbar für
Ensemble-Statistik, aber grenzwertig für präzise Einzelkaskaden-Analyse.

\paragraph{Tesla T4 zum Vergleich.}
Auf der Tesla T4 läuft dieselbe Konfiguration in 19{,}09\,s -- kaum schneller
als die 64 EPYC-Kerne (22{,}5\,s). Der Grund ist die stark beschnittene
fp64-Rechenleistung der T4 (Turing, sm\_75): $\approx 0{,}25$\,TFLOPS
(1/32 der fp32-Leistung) gegenüber 9{,}7\,TFLOPS fp64 des A100.

\begin{center}\small
\begin{tabular}{@{}lrrr@{}}
\toprule
System & $T$ [s] & Speedup vs.\ 1 Kern & Speedup vs.\ 64 Kerne \\
\midrule
CPU 1 Kern    & 1152{,}9 & 1{,}0$\times$ & -- \\
CPU 64 Kerne  &   22{,}5 & 51$\times$    & 1{,}0$\times$ \\
Tesla T4      &   19{,}1 & 60$\times$    & 1{,}2$\times$ \\
A100 fp64     &    1{,}58 & 731$\times$  & 14$\times$ \\
A100 fp32     &    0{,}999 & 1154$\times$ & 23$\times$ \\
\bottomrule
\end{tabular}
\end{center}

\FloatBarrier
% =====================================================================
\section{GPU-Sättigung: Skalierung mit der Problemgröße}
\label{sec:saturation}

\begin{figure}[htbp]
  \centering
  \includegraphics[width=0.85\linewidth]{gpu_saturation.png}
  \caption{GPU-Sättigungskurve: Durchsatz in Matom/s als Funktion der
    Systemgröße $N$ (A100, fp64). Bereits ab $N \approx 250\,000$ Atomen
    ist der Durchsatz konstant auf $\approx 3800$\,Matom/s -- die GPU hat
    genug Threads um alle SMs vollständig auszulasten.
    Skript: \texttt{scripts/python/plot\_gpu\_saturation.py}.}
  \label{fig:saturation}
\end{figure}

Die Plateaustruktur belegt, dass die GPU ab $N \approx 250$\,k voll
gesättigt ist. Für $N < 250$\,k reichen die Threads nicht aus, um alle
108 SMs des A100 zu füllen. Für $N > 250$\,k skaliert die Laufzeit
linear mit $N$ -- der Durchsatz ist konstant.

\FloatBarrier

\newpage
% =====================================================================
\section{ncu-Profiling: Warum ist der A100 trotzdem kein Optimum?}
\label{sec:ncu}

\subsection{SpeedOfLight: Compute vs.\ Memory}

NVIDIA Nsight Compute (ncu) misst den Kernel als Prozentsatz des
theoretischen Peak-Durchsatzes:

\begin{figure}[htbp]
  \centering
  \includegraphics[width=0.9\linewidth]{ncu_profiling.png}
  \caption{ncu-Profiling des \texttt{k\_force}-Kernels.
    Links: Kernel-Zeitanteile pro Schritt (CPU vs.\ A100 fp64).
    Mitte: SpeedOfLight-Throughput -- T4 ist compute-bound (87\,\%),
    A100 ist latenz-limitiert (45\,\% Compute, 37\,\% Memory).
    Rechts: Theoretische Okkupanz durch Registerlimitierung.
    Skript: \texttt{scripts/python/plot\_ncu.py}.}
  \label{fig:ncu}
\end{figure}

\paragraph{Befund A100.}
\begin{itemize}[itemsep=2pt]
\item Compute-Throughput: \textbf{45\,\%} (nicht compute-bound).
\item Memory-Throughput: \textbf{37\,\%} (nicht vollständig memory-bound).
\item Schlussfolgerung: Der Kernel ist \emph{latenz-limitiert} -- Threads
  warten auf Daten aus dem DRAM, können nicht rechnen, und der Scheduler
  hat nicht genug alternative Warps um die Latenz zu verstecken.
\end{itemize}

\paragraph{Befund Tesla T4.}
\begin{itemize}[itemsep=2pt]
\item Compute-Throughput: \textbf{87\,\%} -- klar compute-bound.
\item Memory-Throughput: \textbf{8\,\%} -- sehr gering.
\end{itemize}

Der Unterschied: Die T4 hat viel weniger fp64-Rechenleistung
($\approx 0{,}25$\,TFLOPS vs.\ 9{,}7\,TFLOPS beim A100), aber dieselbe
Kernel-Struktur. Die begrenzte
Rechenkapazität der T4 wird voll ausgelastet; beim A100 ist die Bandbreite
der Engpass -- aber nur zu 37\,\% ausgenutzt.

\subsection{Okkupanz: Registerlimitierung}
\label{sec:occupancy}

\infobox{Okkupanz-Rechnung (fp64, A100)}{%
  Register pro Thread: 62 (gemessen mit \texttt{nvcc --ptxas-info}).
  Register-Budget pro SM: 65\,536.
  Maximale Blöcke/SM: $\lfloor 65536 / (62 \times 256) \rfloor = 4$.
  Warps/SM bei 4 Blöcken: $4 \times 8 = 32$.
  Max Warps/SM (A100): 64. $\Rightarrow$ \textbf{Okkupanz = 32/64 = 50\,\%}.
}

Mit nur 50\,\% Okkupanz hat der Scheduler nur halb so viele Warps
verfügbar, um Speicherlatenz zu verstecken. Bei fp32 sinkt der
Registerbedarf auf $\approx 48$ Register/Thread, was 5 Blöcke/SM
ermöglicht und die Okkupanz auf 62\,\% erhöht.

\paragraph{Naheliegender Hebel.}
Um die Okkupanz zu verbessern, müsste die Register-Allokation reduziert
werden, damit mehr Warps gleichzeitig je SM laufen und die Speicherlatenz
besser versteckt wird. Der direkteste Weg dazu ist das Pragma
\texttt{\_\_launch\_bounds\_\_(maxThreads, minBlocksPerSM)}, das \texttt{ptxas}
zwingt, die Register je Thread so weit zu begrenzen, dass \texttt{minBlocksPerSM}
Blöcke gleichzeitig Platz finden. Ob dieser Hebel \emph{tatsächlich} Leistung
bringt, lässt sich direkt messen -- der nächste Abschnitt tut das.

\FloatBarrier
\subsection{Experiment: Occupancy-Hebel via \texttt{\_\_launch\_bounds\_\_}}
\label{sec:launchbounds}

Vier Builds wurden bei sonst identischem Code verglichen ($1414 \times 1414
\approx 2$\,Mio Atome, fp64, Blockgröße 256, A100): \textbf{base} ohne
Begrenzung sowie \textbf{lb4/lb6/lb8} mit \texttt{minBlocksPerSM} $= 4/6/8$.
Die Absolutwerte dieser Messreihe (base: 3435\,Matom/s) stammen aus einer
separaten Kampagne und liegen $\approx 10\,\%$ unter dem Headline-Benchmark
aus Abschnitt~\ref{sec:a100bench}; für die Aussage zählt der Vergleich der
Varianten untereinander.

\begin{figure}[htbp]
  \centering
  \includegraphics[width=0.85\linewidth]{launch_bounds_plot.png}
  \caption{Occupancy-Hebel via \texttt{\_\_launch\_bounds\_\_}. Blau: per ncu
    gemessene Okkupanz; orange: Durchsatz relativ zur \texttt{base}-Variante.
    \texttt{lb4} ist nicht separat dargestellt, da es (siehe Tabelle unten)
    mit 62 Reg/Thread identisch zu \texttt{base} ist und keine eigene Wirkung zeigt.
    Die Begrenzung hebt die Okkupanz von 45\,\% (base, 62 Reg) auf 66\,\%
    (lb6, 40 Reg) -- der \textbf{Durchsatz bleibt jedoch flach}
    (3435 $\to$ 3406\,Matom/s). Bei lb8 (32 Reg, theoretisch volle Okkupanz)
    kippt es: \textbf{Register-Spilling} kostet 17\,\% Durchsatz.
    Skript: \texttt{scripts/python/plot\_launch\_bounds.py}.}
  \label{fig:launchbounds}
\end{figure}

\begin{center}\small
\begin{tabular}{@{}lrrrr@{}}
\toprule
Variante & Reg/Thread & Okkupanz & force-Kernel & Durchsatz \\
\midrule
base & 62 & 45{,}1\,\% (ncu) & 276\,µs & 3435\,Matom/s \\
lb4  & 62 & = base          & = base  & 3435\,Matom/s \\
lb6  & 40 & 66{,}4\,\% (ncu) & 262\,µs & 3406\,Matom/s \\
lb8  & 32 & $\sim$100\,\% (theor.) & 387\,µs & 2850\,Matom/s \\
\bottomrule
\end{tabular}
\end{center}

\paragraph{Interpretation.}
\begin{itemize}[itemsep=2pt]
\item \textbf{lb4 ist wirkungslos:} Der Kernel braucht ohnehin nur 62 Register,
  was bereits 4 Blöcke/SM erlaubt -- die Begrenzung ändert nichts.
\item \textbf{lb6 ist der Kernbefund:} Die Register sinken von 62 auf 40, die
  \emph{gemessene} Okkupanz steigt von 45\,\% auf 66\,\% -- der Hebel greift
  also nachweislich. Trotzdem bleibt der Durchsatz im Messrauschen unverändert.
  \emph{Mehr Okkupanz übersetzt sich nicht in mehr Leistung.}
\item \textbf{lb8 zeigt die Kehrseite:} Bei 32 Registern tritt Register-Spilling
  auf, der Kraft-Kernel wird $\approx 40\,\%$ langsamer (Gesamtdurchsatz $-17\,\%$)
  -- der klassische Trade-off zu aggressiver Registerbegrenzung.
\end{itemize}

\infobox{Diagnose experimentell bestätigt}{%
  Die Okkupanz ließ sich gezielt von 45\,\% auf 66\,\% anheben, ohne dass der
  Durchsatz folgte. Das \textbf{bestätigt die ncu-Diagnose}: Der Kraft-Kernel
  ist latenz- bzw.\ bandbreitengebunden, \emph{nicht} occupancy-limitiert.
  Zusätzliche Okkupanz ist erzeugbar, aber wirkungslos -- der Engpass liegt
  am Datenlayout (SoA-Bandbreite, Zugriffskoaleszenz), nicht an der Zahl
  aktiver Warps. Ein vollwertiges Negativergebnis: der laut Diagnose richtige
  Hebel wurde umgesetzt, vermessen und als wirkungslos nachgewiesen.
}

\FloatBarrier
\subsection{Weitere Optimierungsversuche: Launch-Overhead und Präzision}
\label{sec:cudaopt}

Die ncu-Diagnose (Abschnitt~\ref{sec:ncu}) nennt zwei naheliegende Hebel
neben der Okkupanz: den Host-seitigen Kernel-Start-Overhead und die
Rechenpräzision. Beide wurden zusätzlich zum Okkupanz-Experiment aus
Abschnitt~\ref{sec:launchbounds} gezielt getestet und gemeinsam mit diesem
in einer konsolidierten Messreihe verglichen (vier Varianten, sonst
identischer Code, $\approx 2$\,Mio Atome, fp64).

\begin{figure}[htbp]
  \centering
  \includegraphics[width=0.95\linewidth]{cuda_opt_plot.png}
  \caption{Vier GPU-Varianten im Vergleich: \texttt{base}, \texttt{lb}
    (\texttt{\_\_launch\_bounds\_\_}, 40 statt 62 Register), \texttt{graph}
    (CUDA Graphs) und \texttt{lb+graph}. Links: die Okkupanz lässt sich
    durch \texttt{lb} von 45\,\% auf 66\,\% anheben. Rechts: der
    Gesamtdurchsatz bleibt über alle Varianten hinweg flach bei
    $\approx 3800$\,Matom/s ($\pm 1\,\%$, im Bereich der Messstreuung).
    Skript: \texttt{scripts/python/plot\_cuda\_opt.py},
    Daten: \texttt{results/cuda\_opt\_results.csv}.}
  \label{fig:cudaopt}
\end{figure}

\paragraph{Launch-Overhead (CUDA Graphs).}
Pro Zeitschritt werden fünf Kernel gestartet; über tausende Schritte
summiert sich der Host-seitige Start-Overhead. Mit CUDA Graphs wird die
feste Kernel-Sequenz einmalig aufgezeichnet (\texttt{cudaStreamBeginCapture})
und anschließend pro Schritt mit einem einzigen \texttt{cudaGraphLaunch}
abgespielt -- die Physik bleibt dabei unverändert. Der resultierende Gewinn
liegt bei $\approx 1\,\%$ und ist damit vernachlässigbar. Der Energie-Drift
ist in allen vier Varianten identisch.

\paragraph{Präzision (fp32 statt fp64).}
Der Wechsel auf einfache Genauigkeit (Abschnitt~\ref{sec:a100bench}) bringt
Faktor $1{,}58\times$ statt der theoretisch möglichen $2\times$. Der Gewinn
stammt im Wesentlichen aus der halbierten Datenmenge, nicht aus mehr
Rechenleistung -- ein weiterer Beleg für die Bandbreiten- bzw.\ Latenzbindung.
Für die Kaskaden-Ensembles (Teil~II) bleibt es dennoch bei fp64:
Bindungsbrüche hängen an feinen Abstandsdifferenzen (Bruchschwelle
$1{,}15\,L_0$), und der Energie-Drift verschlechtert sich unter fp32
deutlich (0{,}94\,\% statt 0{,}22\,\%).

\infobox{Fazit: drei Hebel, ein Engpass}{%
  Höhere Okkupanz, reduzierter Launch-Overhead und reduzierte Präzision
  verbessern die Leistung jeweils nicht wesentlich und weisen konsistent auf
  denselben Engpass hin: die Speicher-Latenz und -Bandbreite des
  \texttt{k\_force}-Kernels. Eine weiterführende Optimierung müsste am
  Datenlayout und an der Zugriffskoaleszenz ansetzen statt an der
  Auslastung. Damit bestätigt die GPU-Analyse denselben limitierenden
  Faktor, der bereits die CPU- und Skalierungsergebnisse dieser Arbeit
  bestimmt.
}

\FloatBarrier   % GPU-Grafiken (Saettigung, ncu, launch_bounds, cuda_opt) vor dem Fazit platzieren

% =====================================================================
\section{Fazit Teil I}
\label{sec:fazit-hpc}

\paragraph{Der durchgehende rote Faden: Speicherbandbreite.}

\begin{itemize}[itemsep=2pt]
\item \textbf{Einzelkern (Compiler):} \texttt{-O3} = $2\times$. \texttt{-Ofast}
  hilft nicht (und verletzt Korrektheit).
\item \textbf{Strong Scaling:} U-förmig. Superlinear bei 128 Kernen durch
  Cache-Effekt (Datensatz pro Rank unterschreitet die 4-MB-L3-Schwelle).
\item \textbf{Weak Scaling:} Bricht ab 32 Kernen ein -- ges\"attigte Knotenbandbreite.
\item \textbf{Optimierungen:} Non-blocking Halo und Re-Sorting bringen $+0{,}8\,\%$ --
  die Kommunikation ist nicht der Engpass.
\item \textbf{GPU:} A100 ist $731\times$ schneller als 1 Kern, läuft aber nur
  mit 45\,\% Compute- und 37\,\% Memory-Auslastung -- latenz-limitiert.
\item \textbf{Occupancy-Hebel (\texttt{\_\_launch\_bounds\_\_}):} Okkupanz
  gezielt auf 66\,\% angehoben -- Durchsatz bleibt flach. Bestätigt die
  Bandbreiten-/Latenz-Diagnose experimentell statt nur per Profiler.
\item \textbf{CUDA Graphs \& Präzision:} Reduzierter Launch-Overhead
  ($+1\,\%$) und fp32 statt fp64 ($1{,}58\times$ statt theor.\ $2\times$)
  bringen beide keinen wesentlichen Zusatzgewinn -- drei unabhängige Hebel,
  ein Engpass.
\end{itemize}

% =====================================================================
% TEIL II: POWER LAW
% =====================================================================
\newpage
\FloatBarrier
\part{Power Law -- Skalenfreie Defektverteilung}
\label{part:powerlaw}

% =====================================================================
\FloatBarrier
\section{Theoretischer Hintergrund}
\label{sec:powerlaw-theory}

\subsection{Was ist ein Potenzgesetz?}

Eine Wahrscheinlichkeitsverteilung $p(n)$ folgt einem Potenzgesetz, wenn
\begin{equation}
p(n) \propto n^{-(\alpha + 1)}, \qquad n \geq n_\text{min} > 0,
\label{eq:powerlaw}
\end{equation}
mit einem Exponenten $\alpha > 0$. In der Literatur zur Strahlenschaden-Physik
wird die Konvention $F(n) = A\,n^{-S}$ verwendet (Häufigkeit statt
Wahrscheinlichkeitsdichte), mit $S = \alpha + 1$. Charakteristisch für
Potenzgesetze ist das Fehlen einer ausgezeichneten Skala: es gibt keine
``typische'' Clustergröße. Das ist das Kennzeichen \emph{selbstorganisierter
Kritikalität} -- verwandte Phänomene sind Erdbebengrößen, Waldbrände,
Netzwerk-Knotengrade.

\paragraph{CCDF -- Complementary Cumulative Distribution Function.}
Direkt aus den Daten schätzbar ist die \emph{komplementäre kumulative
Verteilungsfunktion}
\begin{equation}
P(N \geq n) = \int_n^\infty p(n')\,dn' \propto n^{-S+1},
\label{eq:ccdf}
\end{equation}
die im log-log-Plot als Gerade mit Steigung $-(S-1)$ erscheint.
Die CCDF hat gegenüber einem Histogramm den Vorteil, keine Binning-Entscheidungen
zu benötigen und den Schwanz der Verteilung weniger zu glätten.

\subsection{Sand et al.\ (2013): Potenzgesetz in 3D-Wolfram}
\label{sec:sand}

Sand et al.\ führten vollständige MD-Simulationen von 150-keV-Kollisionskaskaden
in Wolfram durch (PARCAS-Code, $\approx 100$ Kaskaden, drei verschiedene
interatomare Potenziale). Wolfram (W) ist das Kandidatmaterial für die
erste Wand in Fusionsreaktoren wegen hoher Schmelztemperatur,
guter Wärmeleitfähigkeit und geringer Neutronenaktivierung.

\paragraph{Schlüsselbefunde.}
\begin{itemize}[itemsep=2pt]
\item Die Größenverteilung der Zwischengitteratom-Cluster (SIA-Cluster)
  folgt über fast drei Größenordnungen dem Potenzgesetz:
  \begin{equation}
  F(n) = A\,n^{-S}, \quad A = 7{,}45 \pm 1{,}52,\quad S = 1{,}63 \pm 0{,}07
  \label{eq:sand}
  \end{equation}
\item In $\approx 100$ Kaskaden wurde \emph{keine einzige} Subkaskaden-Fragmentierung
  beobachtet. Das Schmelzvolumen bleibt kompakt (zusammenhängend, $\sim 30\,000$
  Atome, mittlere Energiedichte $\approx 4$\,eV/Atom).
\item \textbf{Kontrast zu Eisen:} In Fe fragmentieren Kaskaden bereits ab
  $\approx 50$\,keV. W-Atome sind schwerer, Kollisionsquerschnitte größer,
  Rückstoßatome langsamer -- deshalb kompaktere Schmelzbereiche.
\item Die ``fraktale Natur'' der Kaskade (selbstähnliche Verzweigungsstruktur)
  wurde direkt visualisiert. Sie ist die geometrische Ursache des Potenzgesetzes.
\end{itemize}

\paragraph{Versetzungsschleifen-Charaktere.}
Große Schleifen ($> 80$ SIAs) haben durchgehend
$\bm{b} = \frac{1}{2}\langle111\rangle$; kleinere zeigen zu etwa gleichen
Teilen beide Burgers-Vektoren ($\frac{1}{2}\langle111\rangle$ und
$\langle100\rangle$). Die Wahrscheinlichkeit einer Schleife mit $> 150$
Defekten beträgt ca.\ 3\,\% -- in guter Übereinstimmung mit Experimenten.

\subsection{De Backer et al.\ (2016): Subkaskaden und das fragmentierte Regime}
\label{sec:debacker}

De Backer et al.\ erweitern die Analyse auf hohe PKA-Energien, bei denen
die Kaskade in räumlich getrennte \emph{Subkaskaden} zerfällt. Sie verwenden
eine Subdomain-Dekompositionsmethode kombiniert mit der Binärkollisions-Näherung.

\paragraph{Fragmentierungsenergie $E_\text{fr}$.}
In allen 20 untersuchten Metallen (u.a.\ W, Fe, Be, Zr) existiert eine
charakteristische Energie $E_\text{fr}$, oberhalb derer Subkaskaden entstehen.
Die mittlere Anzahl der Sub-Domänen wächst als:
\begin{equation}
\nu(E) = \begin{cases}
1 & E \leq 2^{-1/P} E_\text{fr} \\
2\,\bigl(E/E_\text{fr}\bigr)^P & E \geq 2^{-1/P} E_\text{fr}
\end{cases}
\label{eq:subkaskaden}
\end{equation}
mit $P \in [0{,}8;\,1{,}2]$. Für Wolfram liegt $E_\text{fr} \approx 160$\,keV;
über alle untersuchten Metalle reicht die Schwelle von 2 bis
$\approx 500$\,keV. Die 2D-Kaskaden dieser Arbeit zeigen die Morphologie des
nicht-fragmentierten Regimes -- ein zusammenhängendes Schadensgebiet statt
räumlich getrennter Subkaskaden (Abb.~\ref{fig:damage}) -- und liegen damit
unterhalb der (modellabhängigen) Fragmentierungsschwelle, konsistent mit dem
sauberen Potenzgesetz, das hier beobachtet wird.

\paragraph{Generalisiertes Skalierungsgesetz.}
Oberhalb von $E_\text{fr}$ ergibt sich ein modifiziertes Skalierungsgesetz:
\begin{equation}
G(n, E) \sim n^{-S}\!\left[1 - \left(\frac{n}{n_\text{fr}}\right)^{1+S-\frac{3}{2}Q(E)}\right],
\qquad E > E_\text{fr}
\label{eq:debacker}
\end{equation}
Der Korrekturterm in der eckigen Klammer dämpft große Cluster: oberhalb
von $E_\text{fr}$ ist die maximale Subkaskaden-Größe durch $E_\text{fr}$
begrenzt, unabhängig von der Gesamtenergie. Dies erklärt die kontraintuitive
Beobachtung, dass \emph{hochenergetische Kaskaden weniger große Defektcluster
erzeugen als niederenergetische}.

\paragraph{Verbindung zur eigenen Simulation.}
Die 2D-Ergebnisse (Abschnitt~\ref{sec:powerlaw-results}) zeigen dasselbe
skalenfreie Verhalten wie Sand et al. -- allerdings im nicht-fragmentierten
Regime, mit niedrigerem Exponenten (2D vs.\ 3D erwartet eine Verschiebung).
Für höhere $E_\text{PKA}$-Werte wäre analog zu De Backer ein
Finite-Size-Cutoff durch Gittergrenzen zu erwarten.

% =====================================================================
\newpage
\section{Ensemble-Setup}
\label{sec:ensemble-setup}

\paragraph{Parameter.}
\begin{itemize}[itemsep=2pt]
\item 1000 unabhängige Kaskaden: serieller Code, 128 Instanzen gleichzeitig
  (\texttt{xargs -P 128}).
\item Variiert: Einschlagort (zufällig), Einschlagwinkel (uniform in $[0, 2\pi)$),
  PKA-Energie (log-uniform in $[200; 2500]$).
\item Gitter: $250 \times 250$ Atome, $dt = 3 \times 10^{-4}$, 12\,000 Schritte.
\item Cluster-Erkennung: Union-Find mit Schwellwert-Verschiebung $> 0{,}5\,L_0$
  und Verbindungsradius $< 1{,}5\,L_0$.
\end{itemize}

\paragraph{Zwei Varianten: Heal vs.\ No Heal.}
\begin{itemize}[itemsep=2pt]
\item \textbf{Heal:} Bindungen können sich nach dem Bruch wieder schließen,
  falls die Atome sich re-annähern.
\item \textbf{No Heal:} Einmal gebrochene Bindungen bleiben permanent offen.
\end{itemize}

Ergebnis: 1000 Läufe mit Healing ergeben \textbf{2391 gepoolte Cluster};
1000 Läufe ohne Healing ergeben \textbf{2335 gepoolte Cluster}. Die Zahl
der Cluster ist fast identisch -- das Healing beeinflusst die Statistik kaum,
weil die Fragmentierung in der schnellen ballistischen Phase festgelegt wird,
bevor das Healing-Fenster greift. Abbildung~\ref{fig:damage} zeigt die
Schadensmorphologie beider Varianten im direkten Vergleich.

\begin{figure}[htbp]
  \centering
  \includegraphics[width=0.9\linewidth]{vid_S2_damage.png}
  \caption{Schadenskarte einer Kaskade (Zoom $\pm 15\,L_0$ um den Einschlag,
    Lauf 999 aus dem Ensemble). Farbe = Verschiebung $\log|u|$; rote Segmente
    = Bindungen, deren aktueller Abstand die Bruchschwelle $1{,}15\,L_0$
    überschreitet (aus den Endzustands-CSVs geometrisch rekonstruiert, da
    keine separate Bindungsliste vorliegt). Links: ohne Healing (18 gerissene
    Bindungen im Zoomfenster). Rechts: mit Healing (11). Beide zeigen die
    typische Morphologie: dichter Kern plus verzweigte Ausläufer. Healing dünnt
    den Schaden etwas aus, ändert aber die Skalierungsstruktur nicht.
    Skript: \texttt{scripts/python/plot\_damage\_map.py}.}
  \label{fig:damage}
\end{figure}

% =====================================================================
\FloatBarrier
\section{Fit-Methodik: Maximum-Likelihood-Estimation}
\label{sec:powerlaw-method}

\subsection{Clauset et al.\ -- MLE statt Histogramm-Fit}

Ein linearer Fit im log-log-Histogramm ist \emph{kein} statistisch valider
Schätzer für den Potenzgesetz-Exponenten (Bias durch Binning, falsche
Fehlerstruktur). Stattdessen kommt Maximum-Likelihood-Estimation (MLE)
nach Clauset et al.\ (2009) zum Einsatz. Für ganzzahlige Daten mit unterem Cutoff $n_\text{min}$
ist der MLE-Schätzer:
\begin{equation}
\hat{S} = 1 + n_\text{Stichprobe}\left[\sum_{i=1}^{n_\text{Stichprobe}}
\ln\frac{x_i}{x_\text{min} - 1/2}\right]^{-1},
\label{eq:mle}
\end{equation}
kombiniert mit dem Kolmogorov-Smirnov-Test zur Wahl des optimalen $x_\text{min}$.

\subsection{Die $x_\text{min}$-Falle}
\label{sec:xmin-falle}

Naive Minimierung des KS-Abstands kann in einem ersten Auswertungsversuch in
den Verteilungsschwanz ``fliehen'' und einen Artefakt-Exponenten
$S \approx 29{,}6$ bei $x_\text{min} \approx 2571$ liefern (weit außerhalb
des Datenbereichs). Die KS-Minimierung findet dann den Bereich hinter der
``Klippe'' (Finite-Size-Cutoff) als besten Fit -- aber nur für die drei
verbleibenden Punkte dort.

\paragraph{Lösung.}
$x_\text{min}$-Scan nur im datenreichen Rumpf:
\begin{itemize}[itemsep=2pt]
\item Minimum: $\geq 5\,\%$ der Datenpunkte über $x_\text{min}$ (mind.\ 50 Punkte),
\item Maximum: $x_\text{min} < $ 90-Perzentil des Datensatzes.
\end{itemize}
Damit wird der Artefakt ausgeschlossen und der eigentliche Potenzgesetz-Rumpf
robust gefittet.

% =====================================================================
\FloatBarrier
\section{Ergebnisse: Power Law in 2D}
\label{sec:powerlaw-results}

Abbildung~\ref{fig:pooling} zeigt zunächst die statistische Konvergenz:
erst ab $N \sim 50$ gepoolten Läufen wird der Potenzgesetz-Rumpf der CCDF
sichtbar -- die 1000 Läufe des Ensembles sind also nicht Luxus, sondern
Voraussetzung.

\begin{figure}[htbp]
  \centering
  \includegraphics[width=\linewidth]{E1_pooling_progression.png}
  \caption{Pooling-Progression. Von $N = 1$ bis $N = 500$ werden zufällig
    gezogene Läufe gepoolt und die CCDF der Clustergrößen geplottet.
    Bei $N = 1$ liegt genau ein Punkt. Mit wachsendem $N$ füllt sich
    die Verteilung; ab $N \sim 50$ ist die Potenzgesetzkurve sichtbar.
    \emph{Statistisch aussagekräftige Aussagen erfordern ausreichend viele Kaskaden.}}
  \label{fig:pooling}
\end{figure}

\begin{figure}[htbp]
  \centering
  \includegraphics[width=0.62\linewidth]{E2_powerlaw_overlay.png}
  \caption{Power Law -- Overlay Heal vs.\ No Heal (je 1000 Läufe, 2391 bzw.\
    2335 Cluster). CCDF im log-log-Plot. Beide Kurven fallen über fast drei
    Größenordnungen als Gerade: skalenfreies Potenzgesetz. Der MLE-Fit liefert
    $S = 1{,}41$ (Heal) und $S = 1{,}38$ (No Heal) bei $x_\text{min} = 4$.
    Referenzlinie: Sand et al.\ (2013) 3D-Wolfram mit $S = 1{,}63$.}
  \label{fig:powerlaw-overlay}
\end{figure}

\begin{figure}[htbp]
  \centering
  \includegraphics[width=\linewidth]{E3_powerlaw_pdf_ccdf.png}
  \caption{Duale Darstellung: PDF (log-binned Histogramm, links) und CCDF
    (rechts), jeweils im Vergleich mit und ohne Healing. Der Finite-Size-Cutoff
    am oberen Ende (Klippe) ist in beiden Darstellungen sichtbar: das endliche
    Gitter setzt eine obere Grenze für die Clustergröße.}
  \label{fig:pdf-ccdf}
\end{figure}

\paragraph{Schlüsselergebnis.}
\begin{center}\small
\begin{tabular}{@{}lrrrl@{}}
\toprule
Variante & Läufe & Cluster & $\hat{S}$ & $x_\text{min}$ \\
\midrule
Mit Healing    & 1000 & 2391 & $1{,}41 \pm 0{,}03$ & 4 \\
Ohne Healing   & 1000 & 2335 & $1{,}38 \pm 0{,}03$ & 4 \\
Sand et al. 3D & -- & -- & $1{,}63 \pm 0{,}07$ & -- \\
\bottomrule
\end{tabular}
\end{center}

% =====================================================================
\FloatBarrier
\section{Diskussion}
\label{sec:powerlaw-discussion}

\subsection{Warum $S < 1{,}63$?}

Der gemessene Exponent $S \approx 1{,}4$ liegt systematisch unter dem 3D-Wert
$S = 1{,}63$. Drei Effekte tragen dazu bei:
\begin{enumerate}[itemsep=2pt]
\item \textbf{Dimension (2D vs.\ 3D):}
  In 2D hat die Kaskade einen kleineren Phasenraum für Fragmentierung
  und weniger Wege zur Energieverteilung. Das führt tendenziell zu
  mehr großen Clustern (flacherer CCDF, kleineres $S$).
\item \textbf{Modellvereinfachung:}
  Das Harmonic-Feder-Modell ohne reale Potenziale (ZBL, EAM) erfasst
  nicht die korrekte energieabhängige Streuquerschnittsstruktur.
  In realen Systemen hängt die Kaskadenmorphologie stark vom
  interatomaren Potenzial ab.
\item \textbf{Endliche Systemgröße:}
  Das $250 \times 250$-Gitter setzt einen Finite-Size-Cutoff bei
  großen Clustern (sichtbar als Klippe in Abb.~\ref{fig:pdf-ccdf}).
  MLE mit $x_\text{min} = 4$ fittet nur den Rumpf; der Schwanz
  ist nicht für die Schätzung verantwortlich.
\end{enumerate}

\subsection{Robustheit gegen Healing}

Der Exponent $S$ ändert sich durch Healing von $1{,}38$ auf $1{,}41$ --
eine Verschiebung von nur $\approx 2\,\%$. Die Erklärung ist konsistent
mit dem Mechanismus: die Kaskadengeometrie (und damit die Clustergrößen)
wird in der schnellen \emph{ballistischen Phase} festgelegt. Healing
(langsame Annäherung von Atompaaren) greift erst danach und betrifft
hauptsächlich kleine Cluster an der Peripherie. Die Skalierungsstruktur
ist robust.

\subsection{Verbindung zur De Backer-Theorie}

De Backer et al.\ sagen voraus, dass große Cluster oberhalb von $E_\text{fr}$
durch das modifizierte Skalierungsgesetz~\eqref{eq:debacker} unterdrückt werden.
Im hier verwendeten Modell entspricht der Finite-Size-Cutoff einem analogen Effekt:
das Gitter setzt eine maximale Subkaskaden-Größe. Für höhere PKA-Energien
würde das modifizierte Skalierungsgesetz relevant -- das ist die
natürliche Erweiterungsfrage.

\subsection{Ausblick und offene Fragen}

\begin{itemize}[itemsep=2pt]
\item \textbf{Bootstrap-Konfidenzintervalle:} Per \texttt{powerlaw}-Paket
  (Alstott et al., 2014) wäre ein rigoroseres Konfidenzintervall zu bestimmen.
\item \textbf{Subkaskaden-Schwellwert:} Bei höherer $E_\text{PKA}$ sollte
  gemäß De Backer ein Übergang zu einer gedämpften Verteilung sichtbar
  werden.
\item \textbf{3D-Erweiterung:} Das Potenzgesetz in 3D mit realen Potenzialen
  (ZBL + Einbettungsenergie-Modell) zu reproduzieren, wäre der nächste Schritt
  zum direkten Vergleich mit Sand et al.
\item \textbf{GPU-Ensemble:} Mit dem A100 (731-fache Beschleunigung) könnten
  $10\,000$ statt $1\,000$ Kaskaden in derselben Rechenzeit simuliert werden --
  der statistische Fehler des Exponenten würde von $\pm 0{,}03$ auf
  $\pm 0{,}01$ sinken.
\end{itemize}

% =====================================================================
\FloatBarrier
\section{Machbarkeitsstudie: Skalierung des Ensembles}
\label{sec:ensemble-scaling}

Könnte man den Exponenten $S$ durch einen \emph{deutlich größeren} Lauf
schärfer bestimmen -- mehr Kaskaden (engeres Konfidenzintervall) und größere
Gitter (mehr Dekaden im Potenzbereich)? Diese Frage wurde unter einer harten
Randbedingung untersucht: $\approx 2000$ verbleibende CPU-Stunden auf Noctua\,2,
der Lauf musste also im ersten Versuch gelingen. Daraus folgte das Prinzip
\textbf{„erst messen, dann rechnen"} -- jeder Budgetkennwert wurde an einer
kleinen Probe bestimmt, statt ihn zu schätzen. Abbildung~\ref{fig:ensemble}
fasst die beiden entscheidenden Messungen zusammen.

\begin{figure}[htbp]
  \centering
  \includegraphics[width=\linewidth]{ensemble_scaling.png}
  \caption{Voruntersuchung zum skalierten Ensemble-Lauf.
    \textbf{Links:} MPI-Strong-Scaling eines einzelnen $1000\times1000$-Laufs
    (1\,Mio Atome). Die Skalierung ist bei 64 und 128 Kernen \emph{superlinear}
    ($87\times$ bzw.\ $171\times$) -- derselbe Cache-Effekt wie im Hauptteil
    (Abschnitt~\ref{sec:cache}). Daraus folgt die Produktionswahl 128 Kerne/Lauf:
    ein Lauf belegt einen Knoten allein und behält die superlineare Effizienz.
    \textbf{Rechts:} Energie-Kalibrierung. Mit steigender PKA-Energie verschmilzt
    der Schaden zu \emph{einem} perkolierten Cluster (größter Cluster
    $\approx$ alle verschobenen Atome); ab $\approx 15\,000$ überrennt er das
    Gitter. Ein einzelner Großlauf liefert also keine breite Größenverteilung.
    Skript: \texttt{scripts/python/plot\_ensemble\_scaling.py}.}
  \label{fig:ensemble}
\end{figure}

\paragraph{Budget-Probe (seriell).}
Eine voll gepackte Welle von 128 gleichzeitigen Läufen (realistisch, da sie
sich die Knotenbandbreite teilen) zeigt \emph{überlineare} Kosten mit der
Gittergröße ($\approx N_x^{3{,}5}$ statt $N_x^2$) -- wieder die geteilte
Speicherbandbreite als Engpass:
\begin{center}\small
\begin{tabular}{@{}rrrr@{}}
\toprule
Gitter & Atome & Wandzeit 128er-Welle & Bytes/Lauf \\
\midrule
$250\times250$ & 62\,500  & 39{,}97\,s  & 1{,}63\,MB \\
$400\times400$ & 160\,000 & 209{,}07\,s & 4{,}23\,MB \\
\bottomrule
\end{tabular}
\end{center}

\paragraph{MPI-Skalierung bei 1\,Mio Atomen.}
Der linke Plot in Abb.~\ref{fig:ensemble} basiert auf folgenden Wandzeiten;
die Effizienz übersteigt bei 64/128 Kernen 100\,\% (cache-resident):
\begin{center}\small
\begin{tabular}{@{}rrrr@{}}
\toprule
Kerne & Wandzeit & Speedup & Effizienz \\
\midrule
1   & 533{,}56\,s & 1{,}0   & 100\,\% \\
16  & 59{,}43\,s  & 9{,}0   & 56\,\%  \\
32  & 19{,}70\,s  & 27{,}1  & 85\,\%  \\
64  & 6{,}11\,s   & 87{,}3  & 136\,\% \\
128 & 3{,}12\,s   & 171{,}0 & 134\,\% \\
\bottomrule
\end{tabular}
\end{center}

Für den Großlauf wurde die MPI-Version um eine \texttt{MPI\_Gatherv}-Ausgabe
des per-Atom-Endzustands erweitert (bislang gab sie nur aggregierte Größen aus).
Die Validierung bestätigte die Korrektheit: MPI mit 1 vs.\ 4 Rängen ist
\textbf{bit-identisch}; seriell vs.\ MPI weicht erwartungsgemäß pointweise ab
(chaotische Trajektorien, unterschiedliche Gleitkomma-Summationsreihenfolge),
aber die \emph{Verteilung} -- das einzig Relevante für das Power Law -- ist
identisch.

\paragraph{Entscheidung: Lauf nicht durchgeführt.}
Die Infrastruktur war lauffähig und validiert, und das Budget hätte rund 1500
Läufe getragen. Entscheidend dagegen sprach der \emph{wissenschaftliche} Befund
der rechten Bildhälfte: Ein größeres Gitter mit höherer Energie erzeugt nicht
mehr Dekaden, sondern bei hoher Energie einen einzelnen perkolierten Blob. Der
erwartete Erkenntnisgewinn gegenüber dem bestehenden Ergebnis
($S = 1{,}38$/$1{,}41$) rechtfertigte Aufwand und Risiko im gegebenen Rahmen
nicht. Das ist dieselbe Methodenlehre wie bei der GPU-Analyse:
\textbf{erst die richtige Messgröße bestimmen, dann Rechenzeit investieren.}

\clearpage
% =====================================================================
\FloatBarrier
\section{Fazit Teil II}
\label{sec:fazit-powerlaw}

\paragraph{Skalenfreies Verhalten emergiert aus dem Mechanismus.}
Das 2D-Dreiecks\-gitter-Modell -- ohne Quantenmechanik, ohne reale
Wolfram-Potenziale -- reproduziert das skalenfreie Potenzgesetz der
Defektclustergrößen mit einem Exponenten $S \approx 1{,}4$.
Das Ergebnis ist:
\begin{itemize}[itemsep=2pt]
\item \textbf{Robust gegen Healing:} Nur $2\,\%$ Verschiebung.
\item \textbf{Konsistent mit Sand et al.:} Gleiche qualitative Struktur
  (skalenfreies Potenzgesetz), niedrigerer Exponent aus 2D/Modell-Gründen.
\item \textbf{Jenseits des De Backer-Regimes:} Im nicht-fragmentierten
  Bereich, d.h.\ das generalisierte Skalierungsgesetz~\eqref{eq:debacker}
  wird noch nicht aktiv.
\end{itemize}

\paragraph{Methodische Lehre.}
Die $x_\text{min}$-Falle beim KS-basierten Fit -- Artefakt $S \approx 29{,}6$ --
zeigt, dass Potenzgesetz-Fits ohne Daten-Rumpf-Einschränkung täuschen.
MLE nach Clauset et al.\ mit expliziter $x_\text{min}$-Validierung ist die
einzig zuverlässige Methode.

% =====================================================================
\appendix
\FloatBarrier
\section{Anhang: Skripte und Reproduzierbarkeit}
\label{app:scripts}

\paragraph{Python-Plot-Skripte (\texttt{scripts/python/}).}
\begin{center}\small
\begin{tabular}{@{}lll@{}}
\toprule
Skript & Erzeugt & Abbildung \\
\midrule
\texttt{analyze\_scaling.py}    & \texttt{strong\_scaling.png}\footnotemark & \ref{fig:strong}, \ref{fig:phases}, \ref{fig:compute-comm} \\
\texttt{analyze\_weak.py}       & \texttt{weak\_scaling.png}\footnotemark[\value{footnote}]   & \ref{fig:weak} \\
\texttt{analyze\_compiler\_opt.py} & \texttt{compiler\_opt.png}\footnotemark[\value{footnote}] & \ref{fig:compiler} \\
\texttt{analyze\_opt\_compare.py}  & \texttt{opt\_compare.png}\footnotemark[\value{footnote}]  & \ref{fig:opt} \\
\texttt{analyze\_ensemble.py}   & CCDF-Plots                  & \ref{fig:pooling}, \ref{fig:pdf-ccdf} \\
\texttt{plot\_damage\_map.py}   & \texttt{vid\_S2\_damage.png} & \ref{fig:damage} \\
\texttt{overlay\_ccdf.py}       & Overlay-Plot                & \ref{fig:powerlaw-overlay} \\
\texttt{plot\_tts.py}           & \texttt{tts.png}            & \ref{fig:tts} \\
\texttt{plot\_ncu.py}           & \texttt{ncu\_profiling.png} & \ref{fig:ncu} \\
\texttt{plot\_launch\_bounds.py} & \texttt{launch\_bounds\_plot.png} & \ref{fig:launchbounds} \\
\texttt{plot\_gpu\_saturation.py} & \texttt{gpu\_saturation.png} & \ref{fig:saturation} \\
\texttt{plot\_ensemble\_scaling.py} & \texttt{ensemble\_scaling.png} & \ref{fig:ensemble} \\
\texttt{plot\_cache\_argument.py} & \texttt{cache\_argument.png} & \ref{fig:cache} \\
\bottomrule
\end{tabular}
\end{center}
\footnotetext{Im Bericht unter einem umbenannten Dateinamen mit
  \texttt{H}/\texttt{E}/\texttt{M}-Präfix eingebunden
  (z.\,B.\ \texttt{H1\_strong\_scaling.png} statt \texttt{strong\_scaling.png}).
  Diese Version wurde für die Abbildungen im Bericht separat erzeugt und
  entspricht nicht notwendig 1:1 dem aktuellen Stand des hier gelisteten Skripts.}

\paragraph{Aufruf-Beispiele.}
\begin{lstlisting}[basicstyle=\ttfamily\footnotesize, breaklines=true]
# Strong Scaling (1+2 Knoten):
python3 analyze_scaling.py results/mpi_strong_scaling_1node.csv \
                            results/mpi_strong_scaling_2node.csv

# Cache-Argument:
python3 plot_cache_argument.py results/mpi_strong_scaling_1node.csv \
                                results/mpi_strong_scaling_2node.csv

# TTS-Plot (Daten sind hardcoded aus Benchmark-Outputs):
python3 plot_tts.py

# ncu-Profiling (Daten hardcoded aus ncu_report.csv + Benchmark-Outputs):
python3 plot_ncu.py

# Schadenskarte (rekonstruiert Bindungen geometrisch aus den Endzustaenden):
python3 plot_damage_map.py results/ensemble_no_heal_run_0999_state_final.csv \
                            results/ensemble_run_0999_state_final.csv
\end{lstlisting}

\paragraph{Hardware.}
Alle CPU-Messungen auf \emph{Noctua 2} (PC2 Paderborn): EPYC 7763,
128 Kerne/Knoten, 256\,GB RAM, HDR InfiniBand. GPU-Messungen auf demselben
Cluster: NVIDIA A100-SXM4-40GB, CUDA 13.2. T4-Messungen auf einem
separaten System (sm\_75, CUDA 12.x).

\begin{thebibliography}{9}
\bibitem{sand2013}
A.~E.~Sand, S.~L.~Dudarev, K.~Nordlund:
\textit{High-energy collision cascades in tungsten: Dislocation loops
structure and clustering scaling laws}.
EPL (Europhysics Letters) 103 (2013) 46003.

\bibitem{debacker2016}
A.~De Backer, A.~Sand, K.~Nordlund, L.~Luneville, D.~Simeone, S.~Dudarev:
\textit{Subcascade formation and defect cluster size scaling in
high-energy collision events in metals}.
EPL 115 (2016) 26001.

\bibitem{clauset2009}
A.~Clauset, C.~R.~Shalizi, M.~E.~J.~Newman:
\textit{Power-law distributions in empirical data}.
SIAM Review 51(4) (2009), 661--703.
\end{thebibliography}

\end{document}
